\documentclass[11pt,a4paper]{article}

\usepackage[utf8]{inputenc}
\usepackage[T1]{fontenc}
\usepackage[brazilian]{babel}
\usepackage{lmodern}
\usepackage[margin=2.4cm]{geometry}
\usepackage{amsmath,amssymb}
\usepackage{xcolor}
\usepackage[most]{tcolorbox}
\usepackage{enumitem}
\usepackage{hyperref}
\usepackage{fancyhdr}
\usepackage{titlesec}
\usepackage{microtype}
\usepackage{array}
\usepackage{booktabs}

\definecolor{deitelblue}{RGB}{31,73,125}
\definecolor{boapratcolor}{RGB}{0,112,60}
\definecolor{errocolor}{RGB}{178,34,34}
\definecolor{engcolor}{RGB}{31,73,125}
\definecolor{desempcolor}{RGB}{198,89,17}

\hypersetup{colorlinks=true, linkcolor=deitelblue, urlcolor=deitelblue,
  pdftitle={C: Como Programar - Apendice C - Exercicios}}

\pagestyle{fancy}
\fancyhf{}
\fancyhead[L]{\small\textit{C: Como Programar} --- Ap\^endice~C}
\fancyhead[R]{\small Exerc\'icios}
\fancyfoot[C]{\small\thepage}
\renewcommand{\headrulewidth}{0.4pt}

\titleformat{\section}{\Large\bfseries\color{deitelblue}}{\thesection}{1em}{}
\titleformat{\subsection}{\large\bfseries\color{deitelblue}}{\thesubsection}{1em}{}

\newtcolorbox{obseng}{enhanced, breakable, colback=engcolor!8, colframe=engcolor,
  fonttitle=\bfseries, title={\textbf{Observa\c{c}\~ao}},
  left=8pt, right=8pt, top=4pt, bottom=4pt}
\newtcolorbox{boaprat}{enhanced, breakable, colback=boapratcolor!8, colframe=boapratcolor,
  fonttitle=\bfseries, title={\textbf{Dica}},
  left=8pt, right=8pt, top=4pt, bottom=4pt}
\newtcolorbox{resposta}{enhanced, breakable, colback=deitelblue!5, colframe=deitelblue!70,
  boxrule=0.6pt, left=8pt, right=8pt, top=4pt, bottom=4pt}
\newtcolorbox{enunciado}{enhanced, breakable, colback=deitelblue!5, colframe=deitelblue!70,
  boxrule=0.6pt, left=8pt, right=8pt, top=4pt, bottom=4pt}

\title{
  \vspace{-1cm}
  {\Huge\bfseries\color{deitelblue} Ap\^endice C}\\[0.3em]
  {\Large\bfseries Sistemas de Numera\c{c}\~ao}\\[0.8em]
  {\large\itshape Exerc\'icios --- Resolu\c{c}\~ao Did\'atica}
}
\author{}
\date{}

\begin{document}
\maketitle
\thispagestyle{empty}
\vspace{-2em}

\begin{center}
\rule{0.85\textwidth}{0.5pt}\\[0.4em]
\textit{``Resolvemos os 11 exerc\'icios C.16 a C.26. Eles cobrem o mesmo terreno da autorrevis\~ao --- convers\~oes entre bases e representa\c{c}\~ao de negativos --- mas com novos valores e duas curiosidades: a base 12 (C.16) e as bases 6, 13 e 3 (C.17). Cada resposta inclui verifica\c{c}\~ao matem\'atica cruzada.''}\\[0.4em]
\rule{0.85\textwidth}{0.5pt}
\end{center}

\vspace{1em}

% ============================================================
\section{Exerc\'icio C.16 --- Sistema base 12}
% ============================================================

\begin{enunciado}
\textbf{Enunciado.} A base 12 seria mais f\'acil porque 12 \'e divis\'ivel por mais n\'umeros que 10. Qual o menor d\'igito da base 12? Qual o maior? Quais os valores posicionais das 4 posi\c{c}\~oes da direita?
\end{enunciado}

\subsection*{Menor d\'igito}

\begin{resposta}\textbf{0} (zero), assim como em \emph{qualquer} sistema posicional.\end{resposta}

\subsection*{Maior d\'igito}

\begin{resposta}\textbf{Um a menos que a base, ou seja, 11.}\end{resposta}

Como n\~ao temos um \'unico s\'imbolo decimal para 10 e 11, precisamos inventar novos s\'imbolos. Conven\c{c}\~oes hist\'oricas e duozimais propuseram:

\begin{center}
\renewcommand{\arraystretch}{1.2}
\begin{tabular}{|c|c|c|}
\hline
\textbf{Valor} & \textbf{S\'imbolo do livro Deitel} & \textbf{S\'imbolos duozimais hist\'oricos} \\
\hline
10 & A & X (Pitman) ou T (Bell) \\
11 & B & E (Pitman) ou L (Bell) \\
\hline
\end{tabular}
\end{center}

Adotamos as letras A e B no estilo hexadecimal.

\subsection*{Valores posicionais das 4 posi\c{c}\~oes da direita}

S\~ao pot\^encias de 12:

\[
12^0 = \boxed{1}, \quad 12^1 = \boxed{12}, \quad 12^2 = \boxed{144}, \quad 12^3 = \boxed{1728}
\]

\begin{resposta}
\begin{center}
\begin{tabular}{|c|c|c|c|}
\hline
P3 & P2 & P1 & P0 \\
\hline
1728 & 144 & 12 & 1 \\
\hline
\end{tabular}
\end{center}
\end{resposta}

\begin{obseng}
\textbf{Por que duod\'ecimo seria mais conveniente?} 12 tem divisores 1, 2, 3, 4, 6, 12 (6 divisores), enquanto 10 tem apenas 1, 2, 5, 10 (4 divisores). Em base 12, a divis\~ao por 3 e 4 d\'a resultado exato e simples (sem d\'izimas), o que ajudaria muito em c\'alculos cotidianos. O sistema sexagesimal (60), usado pelos sumerios para tempo e \^angulos, leva a ideia ainda mais longe: 60 tem 12 divisores.
\end{obseng}

% ============================================================
\section{Exerc\'icio C.17 --- Tabela de valores posicionais (bases 6, 13, 3)}
% ============================================================

\begin{enunciado}
\textbf{Enunciado.} Complete a tabela de valores posicionais.
\end{enunciado}

\begin{resposta}
\begin{center}
\renewcommand{\arraystretch}{1.3}
\begin{tabular}{|l|c|c|c|c|}
\hline
\textbf{Sistema} & \textbf{P3} & \textbf{P2} & \textbf{P1} & \textbf{P0} \\
\hline
Decimal & 1000 & 100 & 10 & 1 \\
Base 6  & \textbf{216} & \textbf{36}  & 6  & \textbf{1} \\
Base 13 & \textbf{2197} & 169 & \textbf{13} & \textbf{1} \\
Base 3  & 27   & \textbf{9}   & \textbf{3}  & \textbf{1} \\
\hline
\end{tabular}
\end{center}
\end{resposta}

\subsection*{C\'alculos}

\begin{itemize}[leftmargin=*,itemsep=0pt]
  \item \textbf{Base 6:} $6^0=1$, $6^1=6$, $6^2=36$, $6^3=216$.
  \item \textbf{Base 13:} $13^0=1$, $13^1=13$, $13^2=169$, $13^3=2197$.
  \item \textbf{Base 3:} $3^0=1$, $3^1=3$, $3^2=9$, $3^3=27$.
\end{itemize}

% ============================================================
\section{Exerc\'icio C.18 --- Bin\'ario 100101111010 para octal e hex}
% ============================================================

\subsection*{Para octal: agrupe de 3 em 3 (a partir da direita)}

\[
\underbrace{100}_{4} \, \underbrace{101}_{5} \, \underbrace{111}_{7} \, \underbrace{010}_{2}
\]

\begin{resposta}\textbf{Octal: 4572.}\end{resposta}

\subsection*{Para hexadecimal: agrupe de 4 em 4 (a partir da direita)}

\[
\underbrace{1001}_{9} \, \underbrace{0111}_{7} \, \underbrace{1010}_{A=10}
\]

\begin{resposta}\textbf{Hexadecimal: 97A.}\end{resposta}

\subsection*{Verifica\c{c}\~ao}

$100101111010_2 = 2426_{10}$.

Conferindo octal: $4572_8 = 4 \times 512 + 5 \times 64 + 7 \times 8 + 2 = 2048 + 320 + 56 + 2 = 2426$. ✓

Conferindo hex: $\text{97A}_{16} = 9 \times 256 + 7 \times 16 + 10 = 2304 + 112 + 10 = 2426$. ✓

\newpage
% ============================================================
\section{Exerc\'icio C.19 --- Hexadecimal 3A7D em bin\'ario}
% ============================================================

\[
\underbrace{3}_{0011} \, \underbrace{A}_{1010} \, \underbrace{7}_{0111} \, \underbrace{D}_{1101}
\]

\begin{resposta}\textbf{Bin\'ario: 0011 1010 0111 1101.}\end{resposta}

\subsection*{Verifica\c{c}\~ao}

$\text{3A7D}_{16} = 3 \times 4096 + 10 \times 256 + 7 \times 16 + 13 = 12288 + 2560 + 112 + 13 = 14973_{10}$.

% ============================================================
\section{Exerc\'icio C.20 --- Hexadecimal 765F em octal}
% ============================================================

\textbf{Dica do livro:} converta primeiro para bin\'ario.

\subsection*{Etapa 1: hex 765F para bin\'ario}

\[
\underbrace{7}_{0111} \, \underbrace{6}_{0110} \, \underbrace{5}_{0101} \, \underbrace{F}_{1111}
\]

Bin\'ario: \texttt{0111 0110 0101 1111}.

\subsection*{Etapa 2: bin\'ario para octal (agrupando de 3 em 3)}

S\~ao 16 bits ao todo; precisamos completar com zeros \`a esquerda para totalizar um m\'ultiplo de 3 (18 bits, 6 grupos):

\[
\underbrace{000}_{0} \, \underbrace{111}_{7} \, \underbrace{011}_{3} \, \underbrace{001}_{1} \, \underbrace{011}_{3} \, \underbrace{111}_{7}
\]

\begin{resposta}\textbf{Octal: 073137 (ou simplesmente 73137).}\end{resposta}

\subsection*{Verifica\c{c}\~ao}

$\text{765F}_{16} = 7 \times 4096 + 6 \times 256 + 5 \times 16 + 15 = 28672 + 1536 + 80 + 15 = 30303_{10}$.

E $73137_8 = 7 \times 4096 + 3 \times 512 + 1 \times 64 + 3 \times 8 + 7 = 28672 + 1536 + 64 + 24 + 7 = 30303$. ✓

% ============================================================
\section{Exerc\'icio C.21 --- Bin\'ario 1011110 em decimal}
% ============================================================

\begin{center}
\renewcommand{\arraystretch}{1.2}
\begin{tabular}{|c|c|c|c|c|c|c|c|}
\hline
Posi\c{c}\~ao    & 6   & 5  & 4  & 3 & 2 & 1 & 0 \\
\hline
Valor pos. & 64  & 32 & 16 & 8 & 4 & 2 & 1 \\
\hline
D\'igito    & 1   & 0  & 1  & 1 & 1 & 1 & 0 \\
\hline
Produto   & 64  & 0  & 16 & 8 & 4 & 2 & 0 \\
\hline
\end{tabular}
\end{center}

\[
64 + 16 + 8 + 4 + 2 = \boxed{94_{10}}
\]

\begin{resposta}\textbf{Decimal: 94.}\end{resposta}

% ============================================================
\section{Exerc\'icio C.22 --- Octal 426 em decimal}
% ============================================================

\begin{center}
\renewcommand{\arraystretch}{1.2}
\begin{tabular}{|c|c|c|c|}
\hline
Posi\c{c}\~ao    & 2   & 1  & 0 \\
\hline
Valor pos. & 64  & 8  & 1 \\
\hline
D\'igito    & 4   & 2  & 6 \\
\hline
Produto   & 256 & 16 & 6 \\
\hline
\end{tabular}
\end{center}

\[
256 + 16 + 6 = \boxed{278_{10}}
\]

\begin{resposta}\textbf{Decimal: 278.}\end{resposta}

\newpage
% ============================================================
\section{Exerc\'icio C.23 --- Hexadecimal FFFF em decimal}
% ============================================================

\begin{center}
\renewcommand{\arraystretch}{1.2}
\begin{tabular}{|c|c|c|c|c|}
\hline
Posi\c{c}\~ao        & 3      & 2     & 1   & 0 \\
\hline
Valor pos.     & 4096   & 256   & 16  & 1 \\
\hline
D\'igito        & F (15) & F (15)& F (15) & F (15) \\
\hline
Produto       & 61440  & 3840  & 240 & 15 \\
\hline
\end{tabular}
\end{center}

\[
61440 + 3840 + 240 + 15 = \boxed{65535_{10}}
\]

\begin{resposta}\textbf{Decimal: 65535.}\end{resposta}

\begin{obseng}
$65535 = 2^{16} - 1$ \'e o maior inteiro \emph{sem sinal} represent\'avel em 16 bits. \'E o limite superior de \texttt{unsigned short} em muitas implementa\c{c}\~oes. Tamb\'em conhecido como ``portas TCP/IP m\'aximas'' (65535 portas dispon\'iveis). E \'e um divisor de Mersenne: $2^{16} - 1 = 3 \times 5 \times 17 \times 257$.
\end{obseng}

% ============================================================
\section{Exerc\'icio C.24 --- Decimal 299 em bin, oct, hex}
% ============================================================

\subsection*{Bin\'ario}

$299 = 256 + 32 + 8 + 2 + 1$:

\begin{center}
\renewcommand{\arraystretch}{1.2}
\begin{tabular}{|c|c|c|c|c|c|c|c|c|}
\hline
Valor pos. & 256 & 128 & 64 & 32 & 16 & 8 & 4 & 2 & 1 \\
\hline
D\'igito    & 1   & 0   & 0  & 1  & 0  & 1 & 0 & 1 & 1 \\
\hline
\end{tabular}
\end{center}

\begin{resposta}\textbf{Bin\'ario: 100101011.}\end{resposta}

\subsection*{Octal: agrupando o bin\'ario de 3 em 3}

\[
\underbrace{100}_{4} \, \underbrace{101}_{5} \, \underbrace{011}_{3}
\]

\begin{resposta}\textbf{Octal: 453.}\end{resposta}

\subsection*{Hexadecimal: agrupando o bin\'ario de 4 em 4}

(Completamos com 3 zeros \`a esquerda para fechar 12 bits = 3 grupos.)

\[
\underbrace{0001}_{1} \, \underbrace{0010}_{2} \, \underbrace{1011}_{B=11}
\]

\begin{resposta}\textbf{Hexadecimal: 12B.}\end{resposta}

\subsection*{Verifica\c{c}\~ao}

\begin{itemize}[leftmargin=*,itemsep=0pt]
  \item Bin: $256+32+8+2+1 = 299$. ✓
  \item Oct: $4 \times 64 + 5 \times 8 + 3 = 256+40+3 = 299$. ✓
  \item Hex: $1 \times 256 + 2 \times 16 + 11 = 256+32+11 = 299$. ✓
\end{itemize}

\newpage
% ============================================================
\section{Exerc\'icio C.25 --- Complemento de 1 e 2 de 779}
% ============================================================

\subsection*{Etapa 1: 779 em bin\'ario}

$779 = 512 + 256 + 8 + 2 + 1$:

\begin{center}
\renewcommand{\arraystretch}{1.2}
\begin{tabular}{|c|c|c|c|c|c|c|c|c|c|}
\hline
Valor pos. & 512 & 256 & 128 & 64 & 32 & 16 & 8 & 4 & 2 & 1 \\
\hline
D\'igito    & 1   & 1   & 0   & 0  & 0  & 0  & 1 & 0 & 1 & 1 \\
\hline
\end{tabular}
\end{center}

\begin{resposta}\textbf{Bin\'ario: 1100001011} (10 bits).\end{resposta}

\subsection*{Etapa 2: complemento de UM (inverte todos os bits)}

\begin{center}
\begin{tabular}{r|cccccccccc}
779             & 1 & 1 & 0 & 0 & 0 & 0 & 1 & 0 & 1 & 1 \\
\hline
$\sim 779$ (C1) & 0 & 0 & 1 & 1 & 1 & 1 & 0 & 1 & 0 & 0
\end{tabular}
\end{center}

\begin{resposta}\textbf{Complemento de UM: 0011110100.}\end{resposta}

\subsection*{Etapa 3: complemento de DOIS (C1 + 1)}

\[
0011110100 + 1 = 0011110101
\]

\begin{resposta}\textbf{Complemento de DOIS: 0011110101.}\end{resposta}

\subsection*{Verifica\c{c}\~ao}

\begin{center}
\begin{tabular}{r}
$\phantom{+}1100001011$ \\
$+ \, 0011110101$ \\
\hline
$\phantom{+}\mathbf{(1)}\,0000000000$
\end{tabular}
\end{center}

O bit transbordado fora dos 10 bits \'e descartado: resultado $0000000000 = 0$. ✓

% ============================================================
\section{Exerc\'icio C.26 --- Complemento de dois de $-1$ em 32 bits}
% ============================================================

\begin{enunciado}
\textbf{Enunciado.} Mostre o complemento de dois do valor inteiro $-1$ em uma m\'aquina com inteiros de 32 bits.
\end{enunciado}

\subsection*{Etapa 1: $+1$ em bin\'ario de 32 bits}

\[
+1 = \underbrace{0000\,0000\,0000\,0000\,0000\,0000\,0000\,0001}_{32 \text{ bits}}
\]

\subsection*{Etapa 2: complemento de UM (inverte todos os bits)}

\[
\sim 1 = \underbrace{1111\,1111\,1111\,1111\,1111\,1111\,1111\,1110}_{32 \text{ bits}}
\]

\subsection*{Etapa 3: complemento de DOIS (C1 + 1)}

\[
\sim 1 + 1 = \underbrace{1111\,1111\,1111\,1111\,1111\,1111\,1111\,1111}_{32 \text{ bits}}
\]

\begin{resposta}
\textbf{Complemento de dois de $-1$ em 32 bits: 32 bits 1 (todos)}, ou seja:

\[
\text{1111 1111 1111 1111 1111 1111 1111 1111}_2
\]

Em hexadecimal: \textbf{0xFFFFFFFF}. Em decimal sem sinal: \textbf{4294967295} ($= 2^{32} - 1$).
\end{resposta}

\subsection*{Verifica\c{c}\~ao}

\begin{center}
\begin{tabular}{r}
$\phantom{+}0000\,0000\,0000\,0000\,0000\,0000\,0000\,0001$ \quad ($+1$)\\
$+ \, 1111\,1111\,1111\,1111\,1111\,1111\,1111\,1111$ \quad ($-1$ em C2)\\
\hline
$\phantom{+}\mathbf{(1)}\,0000\,0000\,0000\,0000\,0000\,0000\,0000\,0000$
\end{tabular}
\end{center}

O bit transbordado \'e descartado; o resultado em 32 bits \'e zero. ✓ Confirma que estamos somando o oposto.

\begin{obseng}
\textbf{Por que \texttt{(unsigned)(-1) == UINT\_MAX} em C/C++:} quando voc\^e converte $-1$ (\texttt{int}) para \texttt{unsigned}, n\~ao h\'a convers\~ao real --- s\'o reinterpreta\c{c}\~ao dos mesmos bits. E os bits de $-1$ em C2 s\~ao 1's. Logo, \texttt{(unsigned)(-1)} \'e o maior \texttt{unsigned} poss\'ivel:
\begin{itemize}[leftmargin=*,itemsep=0pt]
  \item Em 32 bits: \texttt{UINT\_MAX} $= 4{.}294{.}967{.}295$
  \item Em 64 bits: \texttt{UINT64\_MAX} $= 18{.}446{.}744{.}073{.}709{.}551{.}615$
\end{itemize}
Esse truque \'e \emph{idiomatico} em C: por exemplo, \texttt{string::npos} \'e definido como \texttt{(size\_t)(-1)}.
\end{obseng}

% ============================================================
\section*{Encerramento}
% ============================================================

Resolvemos os 11 exerc\'icios deste ap\^endice. Todas as convers\~oes foram verificadas matematicamente. Sintetizando:

\begin{itemize}[leftmargin=*,itemsep=0pt]
  \item Em qualquer base $b$, d\'igitos v\~ao de 0 a $b-1$; valores posicionais s\~ao $b^0, b^1, b^2, \ldots$
  \item Bin\'ario $\leftrightarrow$ octal: grupos de 3 bits. Bin\'ario $\leftrightarrow$ hex: grupos de 4 bits.
  \item Para outras bases, multiplique d\'igitos por valores posicionais.
  \item Complemento de 2: inverta os bits e some 1.
  \item Em C++: \texttt{cout << oct << n}, \texttt{<< hex << n}, \texttt{<< dec << n} (Cap.~23).
\end{itemize}

Com este ap\^endice, o livro Deitel encerra o material sobre C++. Pr\'oximos ap\^endices (D, E, F) cobrem outros aspectos t\'ecnicos da linguagem.

\vspace{1em}
\begin{center}\rule{0.5\textwidth}{0.4pt}\end{center}

\end{document}
